\section{Full Schema Specifications and Inter-Annotator Audit}
\label{supp:schemas}

% NEW: per-supplement purpose statement at the top
% (MEDIUM #16). Every supplement file gets one of these.
\textit{Purpose of this supplement.}\;
This supplement (i) lists the full label vocabulary for the three
nested schemas, (ii) specifies legal-pair filters and the
heuristic family-to-label projection used in the KB audit, and
(iii) reports the inter-annotator audit of that projection
including per-disagreement details. The audit is a validity check
on the heuristic projection, not a re-annotation pass.

% =====================================================
\subsection{Schema vocabularies}
\label{supp:schemas-vocab}

The three schemas form a strict refinement chain: each adds one
distinction the previous schema did not make. Table~\ref{tab:s3-schemas}
lists all labels.

\begin{table}[H]
\centering
\caption{Label inventory, three schemas.}
\label{tab:s3-schemas}
\small
\begin{tabularx}{\linewidth}{@{}lX@{}}
\toprule
\textbf{Schema} & \textbf{Labels} \\
\midrule
S\textsubscript{family} (4) &
\code{ASSOCIATION\_GENERAL}, \code{DRUG\_DISEASE},
\code{DRUG\_GENE\_REGULATION}, \code{\_\_NEGATIVE\_\_} \\
\addlinespace
S\textsubscript{pair} (8) &
\code{GENE\_DISEASE}, \code{VARIANT\_DISEASE}, \code{DRUG\_DISEASE},
\code{DRUG\_GENE\_REGULATION}, \code{GENE\_GENE\_ASSOC},
\code{DRUG\_VARIANT\_ASSOC}, \code{ASSOCIATION\_GENERAL},
\code{\_\_NEGATIVE\_\_} \\
\addlinespace
S\textsubscript{mech} (13) &
All eight S\textsubscript{pair} labels, plus five DrugProt-derived
mechanism sub-heads: \code{DGR\_INHIBIT}, \code{DGR\_ACTIVATE},
\code{DGR\_METABOLIC}, \code{DGR\_REGULATE},
\code{DGR\_STRUCTURAL}. The carried-over \code{DRUG\_GENE\_REGULATION}
head receives BioRED drug--gene relations only, because BioRED's
drug--gene-relation annotations (Bind, Cotreatment,
Drug\_Interaction; $\sim$23 instances total) are too few to support
mechanism-specific sub-heads. \\
\bottomrule
\end{tabularx}
\end{table}

S\textsubscript{family} is the lower bound: collapsing further would
merge clinical pair-type categories that the downstream KB audit must
discriminate. S\textsubscript{mech} is the upper bound on the present
evaluation corpus, because finer distinctions would require categories
DrugProt does not consistently annotate.

% =====================================================
\subsection{Legal-pair filters and family-to-label projection}
\label{supp:schemas-filters}

A relation candidate is eligible for training and evaluation under a
given schema only if the head-entity and tail-entity types are
covered by that schema's vocabulary. The legal entity-pair tuples
are GENE--DISEASE, VARIANT--DISEASE, DRUG--DISEASE, DRUG--GENE,
GENE--GENE, and DRUG--VARIANT; all other entity-pair tuples
(e.g., VARIANT--GENE, which the three dropped CIViC targets carry)
fall outside the schema vocabulary and are excluded from both
training and evaluation.

The heuristic family-to-label projection used in the KB audit maps
each (entity-pair family, CIViC evidence-type, schema) tuple to an
expected-label set. The rules are: (i) CIViC \emph{Predictive}
evidence on a (drug, gene) pair maps to
\code{DRUG\_GENE\_REGULATION} under S\textsubscript{family} and
S\textsubscript{pair}, and to the set of all DrugProt-derived
mechanism sub-heads (\code{DGR\_INHIBIT}, \code{DGR\_ACTIVATE},
\code{DGR\_METABOLIC}, \code{DGR\_REGULATE},
\code{DGR\_STRUCTURAL}) plus \code{DRUG\_GENE\_REGULATION} under
S\textsubscript{mech} (set-valued KB argmax accuracy credits any of
these as a correct top-1 prediction); (ii) CIViC \emph{Diagnostic}
or \emph{Prognostic} evidence on a (gene, disease) or (variant,
disease) pair maps to the corresponding pair-specific head
(\code{GENE\_DISEASE} or \code{VARIANT\_DISEASE}) under
S\textsubscript{pair} and S\textsubscript{mech}, and to
\code{ASSOCIATION\_GENERAL} under S\textsubscript{family};
(iii) CIViC \emph{Predisposing} evidence on a (gene, disease) pair
maps as in (ii). The full projection table per (family, evidence
type, schema) tuple is released in
\code{knowledge\_grounded\_evidence\_audit/schema\_projection.json}
at the released code revision and is the input the inter-annotator
audit below validates.

% =====================================================
\subsection{Evaluable target population breakdown}
\label{supp:schemas-population}

% Population breakdown supporting §2.1 of the main text.
The 165 raw CIViC-extracted targets reduce to 162 evaluable targets
after dropping three rows whose endpoints (variant--gene) have no
schema-vocabulary covering set under any of the three schemas. The
162 evaluable rows distribute across two entity-pair families
(gene--drug 95.1\%, variant--disease 4.9\%) and three heuristic-
projected family-level labels (Table~\ref{tab:s3-population}).

\begin{table}[H]
\centering
\caption{162 evaluable goldlite targets, by entity-pair family
\(\times\) heuristic-projected family-level label. The three
dropped rows (variant--gene endpoints, unmapped) are listed for
completeness.}
\label{tab:s3-population}
\small
\begin{tabular}{@{}llr@{}}
\toprule
\textbf{Entity-pair family} & \textbf{Heuristic label}
& \textbf{Rows} \\
\midrule
gene--drug         & \code{DRUG\_GENE\_REGULATION}      & 90 \\
gene--drug         & \code{ASSOCIATION\_GENERAL}        & 64 \\
variant--disease   & \code{ASSOCIATION\_GENERAL}        & 8 \\
\midrule
\textbf{Evaluable total} &                              & \textbf{162} \\
\midrule
variant--gene (unmapped, dropped) & ---                 & 3 \\
\midrule
\textbf{Raw CIViC extraction total} &                   & \textbf{165} \\
\bottomrule
\end{tabular}
\end{table}

% =====================================================
\subsection{Inter-annotator audit of the heuristic projection}
\label{supp:iaa}

\paragraph{Sampling.}
We drew a stratified random sample of 30 audit targets from the 162
evaluable rows using \code{random.Random(42)} with stratification on
entity-pair family and a pre-specified minimum of three rows per
stratum. Pure proportional allocation would have placed 28--29 rows
in the gene--drug stratum and 1--2 rows in the variant--disease
stratum; the per-stratum floor lifted variant--disease to 3 and
reduced gene--drug to 27 to preserve \(n=30\). The minority stratum
is therefore over-represented in the audit sample (10\%) relative
to its population share (4.9\%), which is conservative with respect
to the directional-bias finding because it shifts \(\kappa\)'s
weight toward the stratum that contributes less to operational KB
surfacing.

\paragraph{Second-annotator protocol.}
A second independent human curator was not available within the
project's resourcing window; we therefore used a frontier large
language model (Claude Opus 4.7) at \texttt{temperature = 0} with
deterministic decoding as the second annotator, prompted with the
same family-level label definitions used by the heuristic but
blind to the heuristic's output. We treat this as a model-based
validity check on the projection, not as a human-grade
dual-annotation study; the $\kappa$ reported below is best
interpreted as a lower bound on the agreement two trained human
curators would achieve, since the LLM has no biomedical-curator
training and the labelling task is deliberately scoped to
family-level rather than fine-grained relations. The use of an LLM
as a proxy second annotator on coarse text-classification tasks has
precedent in recent NLP methodology
\citep{gilardi2023llmannotator}. The prompt template is in
Figure~\ref{lst:iaa-prompt}; the full prompt, raw LLM responses,
and per-target rationales are released as \code{disagreements.csv}
and \code{audit\_labels.csv} alongside this supplement.

\paragraph{Prompt template.}

\begin{figure}[H]
\caption{Inter-annotator audit prompt template. Variables in
\{curly braces\} are filled per target.}
\label{lst:iaa-prompt}
\small
\begin{verbatim}
You are a biomedical curator labelling drug-gene-variant-disease
assertions extracted from PubMed abstracts. You will be given an
abstract and two pre-identified entities. Your task is to choose
ONE family-level relation label that best describes the
relationship the abstract asserts between the two entities.

Allowed labels (choose exactly one):
- DRUG_DISEASE
- DRUG_GENE_REGULATION
- GENE_DISEASE
- VARIANT_DISEASE
- ASSOCIATION_GENERAL
- __NEGATIVE__

Output STRICT JSON only, no prose, no markdown:
{"label": "<one of the allowed labels>",
 "confidence": <float 0 to 1>,
 "rationale": "<one sentence>"}

---
Abstract:
{abstract_text}

Entity A: "{entity_a_text}" (type: {entity_a_type},
                             span: chars {a_start}-{a_end})
Entity B: "{entity_b_text}" (type: {entity_b_type},
                             span: chars {b_start}-{b_end})

Your label:
\end{verbatim}
\end{figure}

\paragraph{Realised agreement.}
The realised Cohen's \(\kappa\) on family-level expected labels was
\(0.56\) (95\% bootstrap CI [0.32, 0.80]; \(B = 5{,}000\); 7 of 30
disagreements; 0 parse errors). Under the Landis \& Koch taxonomy
\citep{landiskoch1977} this falls in the moderate band; the point
estimate sits below the conventional 0.61 substantial threshold and
the lower CI bound reaches into the fair band.

\paragraph{Disagreement structure.}
The seven disagreements were entirely directional
(Table~\ref{tab:s3-disagreements}): in every disagreeing case the
LLM second annotator chose \code{\_\_NEGATIVE\_\_} while the
heuristic chose a positive relation. Six of seven concerned PubMed
abstracts in which the CIViC-cited drug was not literally named in
the abstract text; the seventh (gefitinib--EGFR, PMID 22370314)
concerned an abstract that names the drug but not the gene side. In
every case, the LLM, working only from the abstract text, returned
\code{\_\_NEGATIVE\_\_} because one side of the CIViC-cited
(drug, gene) pair was missing from the text; the heuristic
inherited the drug--gene relation directly from the CIViC evidence
row irrespective of textual mention. The cited drugs cluster around
EGFR-pathway tyrosine-kinase inhibitors (gefitinib $\times$ 2,
erlotinib $\times$ 1, dacomitinib $\times$ 1, lapatinib $\times$ 1)
and the HSP90 inhibitor tanespimycin ($\times$ 2).

\begin{table}[H]
\centering
\caption{Per-disagreement breakdown of the 7 of 30 targets where
the heuristic and the LLM second annotator disagreed. Six concern
abstracts in which the CIViC-cited drug is not literally named;
one (gefitinib--EGFR, PMID 22370314) concerns an abstract that
names the drug but not the gene. The full per-target audit log
including PMIDs, abstract texts, and LLM rationales is released as
\code{disagreements.csv} alongside this supplement.}
\label{tab:s3-disagreements}
\small
\begin{tabular}{@{}lllr@{}}
\toprule
\textbf{CIViC-cited drug} & \textbf{Heuristic label}
& \textbf{LLM label} & \textbf{Count} \\
\midrule
Tanespimycin                     & \code{DRUG\_GENE\_REGULATION}
                                 & \code{\_\_NEGATIVE\_\_} & 2 \\
Lapatinib                        & \code{DRUG\_GENE\_REGULATION}
                                 & \code{\_\_NEGATIVE\_\_} & 1 \\
Dacomitinib                      & \code{DRUG\_GENE\_REGULATION}
                                 & \code{\_\_NEGATIVE\_\_} & 1 \\
Gefitinib                        & \code{DRUG\_GENE\_REGULATION}
                                 & \code{\_\_NEGATIVE\_\_} & 1 \\
Gefitinib                        & \code{ASSOCIATION\_GENERAL}
                                 & \code{\_\_NEGATIVE\_\_} & 1 \\
Erlotinib                        & \code{ASSOCIATION\_GENERAL}
                                 & \code{\_\_NEGATIVE\_\_} & 1 \\
\midrule
\textbf{Total}                   & & & \textbf{7} \\
\bottomrule
\end{tabular}
\end{table}

% Heuristic-label attribution traced to disagreements.csv:
% DGR -> NEG (5): tanespimycin x2 (GL_0031, GL_0039), lapatinib x1 (GL_0043),
%   dacomitinib x1 (GL_0068), gefitinib x1 (GL_0118)
% AG -> NEG (2): gefitinib x1 (GL_0070), erlotinib x1 (GL_0131)
% total = 7.

\paragraph{Implication for KB argmax accuracy.}
The directional bias above is recall-favouring with respect to the
abstract-only relation discrimination task: the heuristic labels
abstracts as positive whenever the CIViC evidence row asserts a
relation, even when one side of the cited (drug, gene) pair is
missing from the text. KB argmax accuracy therefore rewards
positive prediction on oncology abstracts whose CIViC entry
asserts a relation that the abstract does not literally
substantiate. We do not amend the heuristic;
\S\ref{sec:discussion} of the main text discusses how this
directional bias plausibly compounds with oncology-projected
training to produce the configuration-arm variance asymmetry.

\paragraph{Reproducibility.}
Sampling, prompt rendering, LLM call, and \(\kappa\) computation
are implemented in
\code{knowledge\_grounded\_evidence\_audit/analysis/inter\_annotator\_audit/}
with replayable SLURM entry point
\code{phase\_b\_review\_supplements.sbatch}. The full audit bundle
\code{phase\_b\_review\_supplements\_20260507.tar.gz} (sha256
\code{08f0e38febad09e1e19f0097b9bf9b86dca7d3af9ceb782ea1092e0b2ef43b17})
contains scripts, prompts, raw LLM responses, the disagreements
file, and the \(\kappa\) summary.